\chapter{Thirty-second interview frameworks}
\label{app:interview-frameworks}

This appendix is the operating manual for answering under time pressure. The
philosophy lives in \Cref{app:interview-review}. Here each topic is a reusable
reasoning template: question, thirty-second close, decision tree,
measurements, follow-ups, and mistakes. Memorize the green boxes. Expand only
where the interviewer asks.

\framework{BER increased with stable power}
\label{fw:ber-stable-power}

\fwquestion{Pre-FEC BER rose while average received power held steady. How do
you debug it?}

\execanswer{Power held, so leave the power ledger. First scope the failure
(unit $\rightarrow$ lot $\rightarrow$ vendor $\rightarrow$ fleet). Apply the
power-versus-quality fork: chase eye, bias, wavelength, noise, or receiver.
Therefore I would isolate with a golden swap and a bias sweep, then decide
contain, retune, or ATP change.}

\begin{dectree}
BER up, power stable
  |
Scope: unit / lot / vendor / fleet
  |
Power fork: power held -> quality or Rx
  |
Golden swap Tx vs Rx
  |-- Tx --> eye / bias / wavelength / RIN
  |-- Rx --> TIA / sensitivity / EQ
  |
Decision -> control (ATP / table / lot)
\end{dectree}

\fwwhy{Stable power rules out average launch and connector loss as the primary
cause. The tree forces scope and the fork before any component name so you
spend hours on the ledger that actually moved
(\Cref{sec:interview-power-fork,sec:interview-worked-sensitivity}).}

\fwmeas{Power meter $\rightarrow$ power ledger intact? $\rightarrow$ optical
path or SI.\\
DCA / bias sweep $\rightarrow$ eye or setpoint? $\rightarrow$ retune vs device.\\
Golden swap $\rightarrow$ Tx or Rx? $\rightarrow$ owner.\\
BER waterfall $\rightarrow$ shift or floor? $\rightarrow$ sensitivity vs noise.}

\fwfollow{\begin{itemize}
\item Why measure power before the DCA?
\item Why not start with RIN?
\item What if cool-down recovers the BER?
\end{itemize}}

\fwmistakes{\begin{itemize}
\item Naming the laser first.
\item Calling ``bad eye'' a root cause.
\item Listing instruments without a decision.
\item Stopping before recurrence control.
\end{itemize}}

\fwdeep{Full prose: \Cref{sec:interview-worked-sensitivity,sec:fm-ber-increase}.
Ledgers: \Cref{sec:interview-margin-ledgers}.}

\framework{Received power decreased}
\label{fw:power-down}

\fwquestion{Received optical power dropped and the link is failing. How do you
debug it?}

\execanswer{Power moved, so stay on the power ledger. First scope the failure.
Confirm with an external meter at a named plane, then bisect source enable,
coupling, connectors, MUX loss, and monitor-PD / APC honesty. Therefore I
would contain if lot-correlated, clean or replace the plant if local, and
update ATP or hygiene rules so it does not recur.}

\begin{dectree}
Power down
  |
Scope
  |
External meter @ named plane
  |
Monitor vs external agree?
  |-- NO --> APC / monitor-PD / cal
  |-- YES --> source / coupling / connector / MUX
  |
Population?
  |-- lot/vendor --> contain + FA
  |-- single --> plant or unit repair
  |
Decision + control
\end{dectree}

\fwwhy{A power drop is not automatically a dying laser. Monitor corruption and
dirty plant fake laser failure. The external meter is the cut that protects
you from trusting a lying APC loop.}

\fwmeas{External power meter $\rightarrow$ true power? $\rightarrow$ optical
path.\\
Inspect / ORL $\rightarrow$ plant dirty or reflective? $\rightarrow$ clean /
replace.\\
LIV $\rightarrow$ device changed? $\rightarrow$ replace vs setpoint.\\
Lot query $\rightarrow$ population? $\rightarrow$ contain vs unit fix.}

\fwfollow{\begin{itemize}
\item Why not trust the module's reported Tx power?
\item When do you stop ship versus clean the connector?
\end{itemize}}

\fwmistakes{\begin{itemize}
\item Trusting monitor-PD without an external meter.
\item Treating one dirty connector as a fleet laser problem.
\end{itemize}}

\fwdeep{Power fork and instruments: \Cref{sec:interview-power-fork}. APC and
calibration: \Cref{sec:laser-calibration}.}

\framework{One weak lane}
\label{fw:weak-lane}

\fwquestion{A single lane in a four-lane transceiver has elevated BER while
received power on that lane looks normal. How would you isolate it?}

\execanswer{First scope: one lane, one unit, or a pattern across the lot?
Compare sibling lanes, then optical-versus-electrical swap to split path from
driver or TIA. Therefore I would fix assembly or the array element that owns
the fault, and add the screen that would have caught it.}

\begin{dectree}
One weak lane
  |
Siblings OK?
  |-- NO --> shared supply / thermal / host
  |-- YES --> lane-local
        |
   Opt vs elec swap
        |-- optical --> FAU / coupling / PIC lane
        |-- electrical --> driver / TIA / SerDes
        |
   Lot pattern?
        |
   Decision: rework / screen / supplier
\end{dectree}

\fwwhy{Sibling comparison is the cheapest population cut inside one module.
Optical-versus-electrical swap prevents weeks of laser FA on a TIA lane
(\Cref{sec:interview-worked-lane}).}

\fwmeas{Sibling BER/power $\rightarrow$ shared or local? $\rightarrow$ scope.\\
Lane swap $\rightarrow$ optical or electrical? $\rightarrow$ owner.\\
LIV on weak lane $\rightarrow$ device or coupling? $\rightarrow$ FAU vs die.}

\fwfollow{\begin{itemize}
\item Why compare siblings before opening the module?
\item What lot pattern would make you call the supplier today?
\end{itemize}}

\fwmistakes{\begin{itemize}
\item Rewriting the module spec for one assembly escape.
\item Skipping the electrical path when power looks fine.
\end{itemize}}

\fwdeep{Worked answer: \Cref{sec:interview-worked-lane}. FAU and parallel
optics appear throughout \Cref{ch:failure-modes}.}

\framework{High-temperature failures}
\label{fw:hot-fail}

\fwquestion{BER worsens at high temperature but average power is stable. What
do you do?}

\execanswer{Power held, so leave the power ledger. First scope the failure and
whether cool-down recovers. At the failing temperature, read eye, bias sweep,
wavelength, and control headroom. Therefore I would fix the table or thermal
design and put that loaded corner in the ATP.}

\begin{dectree}
Hot BER, power stable
  |
Scope + cool-down recovers?
  |-- YES --> operating point / table / control
  |-- NO  --> aging / permanent damage
  |
At failing T: eye, bias sweep, OSA, TEC/heater codes
  |
Control ledger exhausted?
  |-- YES --> thermal design / tuning range
  |-- NO  --> cal table / wavelength / Rx
  |
Decision: retune + ATP corner
\end{dectree}

\fwwhy{Temperature failures are often control or calibration, not a dead
laser. Measuring at room temperature misses the spent ledger
(\Cref{sec:interview-worked-hot-ber}).}

\fwmeas{Cool-down test $\rightarrow$ reversible? $\rightarrow$ thermal vs aging.\\
Bias sweep at hot $\rightarrow$ optimum moved? $\rightarrow$ table.\\
Actuator codes $\rightarrow$ control authority left? $\rightarrow$ thermal
design.\\
OSA $\rightarrow$ spectral alignment? $\rightarrow$ lock / filter.}

\fwfollow{\begin{itemize}
\item Why insist on measuring at the failing temperature?
\item How do you say ``TEC maxed'' in product language?
\end{itemize}}

\fwmistakes{\begin{itemize}
\item Debugging only at ambient.
\item Saying TEC current hit max instead of control ledger exhausted.
\end{itemize}}

\fwdeep{\Cref{sec:interview-worked-hot-ber,sec:fm-temperature}. Control
ledger: \Cref{sec:interview-margin-ledgers}.}

\framework{Laser aging versus calibration drift}
\label{fw:aging-vs-cal}

\fwquestion{How do you distinguish laser aging from calibration drift?}

\execanswer{Aging changes the LIV. Drift changes the setpoint on a healthy
LIV. External LIV plus recovery after recalibration separates them. Time
behavior confirms: monotonic climb versus a step after a table or firmware
change. Therefore I would route aging to life/derate/replace and drift to
table control plus an ATP loaded-corner check.}

\begin{dectree}
Symptom (bias up / BER up)
  |
External LIV vs ship data
  |-- LIV moved --> aging --> life / derate / replace
  |-- LIV OK ----> drift path
        |
   Monitor vs external power
        |
   Bias sweep / table reload recovers?
        |-- YES --> calibration / firmware
        |-- NO  --> deeper FA
\end{dectree}

\fwwhy{Telemetry can look identical for aging and drift. Only an external
reference and a recovery test pick the owner
(\Cref{sec:interview-worked-aging-vs-cal}).}

\fwmeas{External LIV $\rightarrow$ device changed? $\rightarrow$ aging vs
setpoint.\\
Monitor vs meter $\rightarrow$ APC honesty? $\rightarrow$ monitor-PD.\\
Table reload $\rightarrow$ recovers? $\rightarrow$ calibration owner.\\
Stress-hours plot $\rightarrow$ monotonic or step? $\rightarrow$ confirms.}

\fwfollow{\begin{itemize}
\item What is the silent failure mode if the monitor lies?
\item Who owns aging versus who owns the table?
\end{itemize}}

\fwmistakes{\begin{itemize}
\item Recalibrating an aged device and calling it fixed.
\item Sending a table bug to the reliability team.
\end{itemize}}

\fwdeep{\Cref{sec:interview-worked-aging-vs-cal,sec:laser-calibration,sec:laser-thermal-vs-aging}.}

\framework{Second-source qualification}
\label{fw:second-source}

\fwquestion{How would you qualify a second laser or PIC supplier?}

\execanswer{Freeze the requirements slice, not the incumbent datasheet.
Characterize distributions across lots and temperature from the customer
view: OMA, RIN at ORL, wavelength, eye, and life with a named HTOL mechanism.
Therefore I would gate open volume on FAIR, ATP correlation, and split RMA
codes, not on a hero sample.}

\begin{dectree}
Second source
  |
Freeze requirements (customer-visible)
  |
Distributions vs means (lots, T)
  |
HTOL + named mechanism
  |
ATP correlation + FAIR
  |
Fleet codes split by supplier
  |
Decision: qualify / hold / reject
\end{dectree}

\fwwhy{The incumbent proves one process can meet the link. The second source
must meet the same external requirements, not copy a datasheet. Customer view
beats internal curiosity early (\Cref{sec:interview-customer-vendor,sec:interview-worked-second-source}).}

\fwmeas{Multi-lot LIV/RIN/SMSR $\rightarrow$ process spread? $\rightarrow$ risk.\\
HTOL $\rightarrow$ life hypothesis? $\rightarrow$ FIT / derate.\\
ATP on both sources $\rightarrow$ escape risk? $\rightarrow$ ship gate.\\
RMA split $\rightarrow$ field falsifies qual? $\rightarrow$ reopen.}

\fwfollow{\begin{itemize}
\item Why not just match the incumbent datasheet?
\item What customer measurements matter if the part is a black box?
\end{itemize}}

\fwmistakes{\begin{itemize}
\item Qualifying a hero sample.
\item Merging supplier RMA codes.
\end{itemize}}

\fwdeep{\Cref{sec:interview-worked-second-source,sec:supplier-exec}.}

\framework{Validation plan for a new transceiver}
\label{fw:validation-plan}

\fwquestion{How would you validate a new optical transmitter from bring-up
through production?}

\execanswer{Validation is staged uncertainty reduction. Walk bring-up,
characterization, margin, interop, environment, reliability, production, and
fleet. Each stage answers a question the previous could not. Therefore I
would refuse any test that answers no new question and watch the ledgers
margin budgeting says will be spent first.}

\begin{dectree}
New transceiver
  |
Bring-up -> Characterization -> Margin
  |
Interop -> Environment -> Reliability
  |
ATP / manufacturing -> Pilot -> Fleet
  |
Each stage: question? uncertainty? decision?
\end{dectree}

\fwwhy{The ladder is Universal Tree~2. Margin budgeting says stresses consume
margin; stages exist to measure remaining margin from the customer-visible
surface first (\Cref{sec:interview-validation-ladder,sec:interview-margin-budget}).}

\fwmeas{Each ladder stage $\rightarrow$ named uncertainty removed
$\rightarrow$ continue / redesign / tighten ATP / stop ship.\\
Margin sweeps $\rightarrow$ which ledger dies first? $\rightarrow$ telemetry
alarms.}

\fwfollow{\begin{itemize}
\item Which stage removes combination risk?
\item How does margin budgeting change what you instrument in the fleet?
\end{itemize}}

\fwmistakes{\begin{itemize}
\item Treating bring-up as population proof.
\item Running HTOL without a named mechanism.
\end{itemize}}

\fwdeep{\Cref{sec:interview-worked-validation,tab:ladder,sec:interview-customer-vendor}.}

\framework{Fleet issue}
\label{fw:fleet}

\fwquestion{Field telemetry shows rising pre-FEC BER on a subset of racks.
How do you triage?}

\execanswer{First scope: unit, lot, vendor, rack, datacenter, or fleet? Ask
trend and change history. Classify performance versus reliability versus
manufacturability before pulling hardware. Therefore I would contain if
growing and supplier-specific, or monitor-only if tiny, flat, and no customer
impact, with an owner on the next control.}

\begin{dectree}
Fleet symptom
  |
Scope ladder
  |
Rate / trend / customer impact
  |-- tiny, flat, no impact --> monitor only
  |-- growing / supplier --> contain now
  |
Bucket: performance / reliability / manufacturability
  |
Decision + owner + telemetry control
\end{dectree}

\fwwhy{Fleet economics and scope pick the owner before FA. Pulling units
without a bucket wastes the only failing state you had
(\Cref{sec:interview-under-uncertainty,sec:fleet-triage}).}

\fwmeas{Telemetry query $\rightarrow$ scope and trend? $\rightarrow$ contain
vs monitor.\\
Lot/date correlation $\rightarrow$ manufacturability? $\rightarrow$ supplier.\\
Golden host in rack $\rightarrow$ environment vs module? $\rightarrow$ owner.}

\fwfollow{\begin{itemize}
\item When is monitor-only the correct Staff answer?
\item What do you do before the mechanism is known?
\end{itemize}}

\fwmistakes{\begin{itemize}
\item Treating one rack as the fleet.
\item Waiting for DPA before containing a growing lot.
\end{itemize}}

\fwdeep{\Cref{sec:interview-under-uncertainty,tab:fleet-triage}.}

\framework{Supplier escape}
\label{fw:supplier-escape}

\fwquestion{A new date code from Supplier~B fails the hot corner at about 3\%.
What do you do?}

\execanswer{Evidence beats an unknown mechanism. Stop shipment of Supplier~B
affected lots, continue Supplier~A, expand the ATP hot-corner screen, open FA
with DPA on fails, and watch RMA daily. Therefore I would contain today and
keep the FA path open rather than wait for SEM before acting.}

\begin{dectree}
Escape signal
  |
Scope: supplier / lot / corner / rate
  |
Mechanism known?
  |-- NO --> still decide
  |
Contain B + continue A
  |
ATP screen + FA/DPA + telemetry
  |
Close when mechanism + control proven
\end{dectree}

\fwwhy{This is Staff judgment under incomplete information. Product deadlines
do not wait for certainty
(\Cref{sec:interview-under-uncertainty}).}

\fwmeas{Lot genealogy $\rightarrow$ exposure? $\rightarrow$ quarantine list.\\
ATP hot corner $\rightarrow$ escape path? $\rightarrow$ production gate.\\
DPA sample $\rightarrow$ mechanism class? $\rightarrow$ permanent fix.}

\fwfollow{\begin{itemize}
\item Why contain before FA completes?
\item What residual risk do you name to leadership?
\end{itemize}}

\fwmistakes{\begin{itemize}
\item Waiting for root cause before stop-ship.
\item Blocking all vendors for one supplier's lot.
\end{itemize}}

\fwdeep{\Cref{sec:interview-under-uncertainty,tab:interview-decisions}.}

\framework{BER floor}
\label{fw:ber-floor}

\fwquestion{Why can a link show a BER floor that more launch power does not
fix?}

\execanswer{A floor means multiplicative noise: SNR stops improving as power
rises. Prioritize RIN under ORL, MPI from reflections, weak PSRR, and
crosstalk. Therefore I would confirm floor versus shift on a waterfall, then
remove the noise source rather than raise OMA.}

\begin{dectree}
BER vs power
  |
Shape?
  |-- parallel shift --> sensitivity
  |-- floor ---------> multiplicative noise
        |
   ORL / RIN / MPI / PSRR / crosstalk
        |
   Fix noise source, not launch
\end{dectree}

\fwwhy{More power cannot beat noise that scales with the signal. The waterfall
shape is the fast cut between sensitivity and floor
(\Cref{sec:interview-worked-ber-floor,sec:interview-waterfall}).}

\fwmeas{Waterfall $\rightarrow$ shift or floor? $\rightarrow$ next hunt.\\
Controlled ORL $\rightarrow$ reflection-driven RIN? $\rightarrow$ plant.\\
Supply noise $\rightarrow$ PSRR path? $\rightarrow$ filter / layout.}

\fwfollow{\begin{itemize}
\item Why not just increase OMA?
\item How do FEC histograms look for MPI bursts?
\end{itemize}}

\fwmistakes{\begin{itemize}
\item Raising launch into a floor.
\item Skipping ORL when RIN is blamed.
\end{itemize}}

\fwdeep{\Cref{sec:interview-worked-ber-floor,sec:fm-ber-floor,sec:rin}.}

\framework{Intermittent failures}
\label{fw:intermittent}

\fwquestion{The link fails intermittently and often passes when retested. How
do you proceed?}

\execanswer{Preserve the failing state and telemetry before you reseat, clean,
or reboot. Scope time and change: dwell, temperature, vibration, firmware,
connector. Prefer burst/FEC histograms and long dwell over a single golden
retest. Therefore I would contain if lot-correlated, tighten dwell/ATP if
escape, and refuse to close an NFF without a reproduction plan.}

\begin{dectree}
Intermittent
  |
Preserve state / telemetry
  |
Triggers: T, time, mate, vibration, FW
  |
Reproduce with dwell / stress
  |
Scope population
  |
Decision: contain / ATP dwell / monitor
\end{dectree}

\fwwhy{Intermittents die when the evidence is destroyed. NFF is often a triage
failure, not a healthy part (\Cref{sec:fleet-triage}).}

\fwmeas{FEC histogram $\rightarrow$ burst vs Gaussian? $\rightarrow$ MPI /
intermittent.\\
Dwell BER $\rightarrow$ reproduces? $\rightarrow$ ATP change.\\
Mate/demate $\rightarrow$ connector? $\rightarrow$ hygiene / replace.}

\fwfollow{\begin{itemize}
\item What do you refuse to do on first touch?
\item How does high NFF change your answer?
\end{itemize}}

\fwmistakes{\begin{itemize}
\item Reseating before capture.
\item Closing NFF without a reproduction plan.
\end{itemize}}

\fwdeep{\Cref{sec:fleet-triage,tab:failure-analysis-checklist}.}

\framework{Production ATP update}
\label{fw:atp-update}

\fwquestion{A failure escaped to the field that ATP did not catch. How do you
update production test?}

\execanswer{Name the escape path and the measurement that would have caught
it at a named plane and corner. Size the new limit from characterization and
repeatability, correlate stations, and set a reaction plan. Therefore I would
ship the ATP change with an owner and a metric, not a hope that operators will
be careful.}

\begin{dectree}
Escape
  |
Which uncertainty ATP missed?
  |
New measurement / corner / limit
  |
Guard band from repeatability
  |
Station correlation
  |
Reaction plan + owner
  |
Ship ATP change
\end{dectree}

\fwwhy{Recurrence control is usually an ATP or telemetry change. A limit
nobody can trace gets waived under ship pressure
(\Cref{sec:interview-worked-validation}).}

\fwmeas{Escape unit replay $\rightarrow$ missed corner? $\rightarrow$ new
test.\\
Gage R\&R $\rightarrow$ repeatability? $\rightarrow$ guard band.\\
Golden units across stations $\rightarrow$ correlation? $\rightarrow$ ship.}

\fwfollow{\begin{itemize}
\item How do you set the guard band?
\item What is the reaction plan when the new limit fails?
\end{itemize}}

\fwmistakes{\begin{itemize}
\item Adding a test with no reaction plan.
\item Limits from one hero unit.
\end{itemize}}

\fwdeep{Production readiness stage in \Cref{sec:interview-validation-ladder}.}

\framework{Telemetry design}
\label{fw:telemetry}

\fwquestion{What would you put in fleet telemetry, and why?}

\execanswer{Log what discriminates hypotheses: per-lane power, bias, pre-FEC
BER and FEC histograms; module temperature and actuator drive; LOS/LOL and
firmware with context. Alarm on trends and disagreements, not only hard
thresholds. Therefore I would instrument the ledgers margin testing said die
first.}

\begin{dectree}
Telemetry purpose: early margin erosion
  |
Per-lane: power, bias, pre-FEC BER
  |
Module: T, TEC/heater, rails, lock
  |
Events + lot/age/rack context
  |
Alarms: trends / disagreements
\end{dectree}

\fwwhy{Telemetry exists to triage without pulling hardware and to catch dying
units before dead ones (\Cref{sec:interview-worked-telemetry}).}

\fwmeas{Bias at constant power $\rightarrow$ aging? $\rightarrow$ reliability
bucket.\\
Power at constant bias $\rightarrow$ optical path? $\rightarrow$ plant.\\
Actuator near rail $\rightarrow$ control margin? $\rightarrow$ thermal
design.}

\fwfollow{\begin{itemize}
\item Why alarm on actuator headroom?
\item What context fields make lot queries possible?
\end{itemize}}

\fwmistakes{\begin{itemize}
\item Logging every register with no hypothesis.
\item Hard thresholds only, no trends.
\end{itemize}}

\fwdeep{\Cref{sec:interview-worked-telemetry,sec:cmis}.}

\framework{Qualification planning}
\label{fw:qual-plan}

\fwquestion{How do you plan qualification for a new IM/DD module?}

\execanswer{Start from customer-visible requirements and a margin budget:
which stresses will spend which ledgers. Map stresses to mechanisms, run HTOL
only with a named activation energy, and prove ATP can catch escapes at rate.
Therefore I would gate ship on remaining margin after environment, interop,
and life, not on a checklist of rituals.}

\begin{dectree}
Qual plan
  |
Customer requirements + margin budget
  |
Stress -> mechanism -> ledger spent
  |
HTOL with named Ea
  |
ATP catches escapes?
  |
Ship gate = remaining margin OK
\end{dectree}

\fwwhy{Margin budgeting and customer view keep qual from becoming museum of
tests. Stress consumes margin; qual measures what remains
(\Cref{sec:interview-margin-budget,sec:interview-customer-vendor}).}

\fwmeas{Margin sweeps $\rightarrow$ cliff distance? $\rightarrow$ derate or
redesign.\\
HTOL $\rightarrow$ wear-out hypothesis? $\rightarrow$ FIT.\\
ATP correlation $\rightarrow$ factory control? $\rightarrow$ volume.}

\fwfollow{\begin{itemize}
\item What does ``remaining margin'' mean in one sentence?
\item When do you ask for Tx-only or Rx-only samples?
\end{itemize}}

\fwmistakes{\begin{itemize}
\item HTOL without mechanism or $E_a$.
\item Qualifying internals the customer never observes.
\end{itemize}}

\fwdeep{\Cref{sec:interview-validation-ladder,sec:gr468,sec:laser-aging}.}

\framework{Unknown failure}
\label{fw:unknown}

\fwquestion{You must make a ship decision this week, but the physical mechanism
is still unknown. What do you do?}

\execanswer{State evidence, confidence weights, and residual risk. Decide with
today's evidence: contain the scoped population, keep a healthy path shipping,
open FA, and add the ATP or telemetry control that would catch the next
escape. Therefore I would not wait for certainty before ownership actions.}

\begin{dectree}
Unknown mechanism
  |
Evidence + scope + rate/trend
  |
Priors: common modes first
  |
Decision today (contain / ship / derate)
  |
FA path + control + owner
  |
Update when mechanism closes
\end{dectree}

\fwwhy{The job is the best decision with today's evidence. Unknown mechanism
is Framework~15 of Staff judgment, not a reason to freeze
(\Cref{sec:interview-under-uncertainty}).}

\fwmeas{Whatever removes the most uncertainty per hour toward the ship
decision: scope query, golden swap, external power, hot-corner ATP sample.\\
DPA later $\rightarrow$ mechanism class? $\rightarrow$ permanent fix.}

\fwfollow{\begin{itemize}
\item What do you tell leadership about residual risk?
\item Which measurement is worth delaying the decision for?
\end{itemize}}

\fwmistakes{\begin{itemize}
\item Freezing all action until SEM.
\item Making a ship call with no control and no owner.
\end{itemize}}

\fwdeep{\Cref{sec:interview-under-uncertainty,sec:interview-decision-trees,tab:interview-decisions}.}

\keyidea{Open the matching framework, deliver the thirty-second box, walk the
tree, end on the decision and the control. Philosophy is in
\Cref{app:interview-review}; this appendix is how you speak it under pressure.}
